\documentclass[11pt]{article}

\usepackage[a4paper,margin=0.72in]{geometry}
\usepackage{amsmath,amssymb}
\usepackage{booktabs}
\usepackage{tabularx}
\usepackage{array}
\usepackage{enumitem}
\usepackage{xcolor}
\usepackage{microtype}
\usepackage{fancyhdr}
\usepackage{listings}

\pagestyle{fancy}
\fancyhf{}
\lhead{DA3410}
\rhead{Viva Question Set-1}
\cfoot{\thepage}

\setlength{\parindent}{0pt}
\setlength{\parskip}{4pt}
\setlist[itemize]{leftmargin=1.4em,itemsep=1pt,topsep=2pt}
\setlist[enumerate]{leftmargin=1.6em,itemsep=1pt,topsep=2pt}

\lstset{
  language=Python,
  basicstyle=\ttfamily\small,
  columns=fullflexible,
  keepspaces=true,
  showstringspaces=false,
  frame=single,
  framerule=0.3pt,
  aboveskip=4pt,
  belowskip=4pt,
  xleftmargin=0.5em,
  xrightmargin=0.5em
}

% Student copy: answers and marking rubrics are intentionally omitted
% from the source, not merely hidden by a toggle.
\newcounter{question}
\newenvironment{question}{%
  \refstepcounter{question}
  \par\medskip
  \noindent\textbf{Q\thequestion.}\par\smallskip
}{\par\medskip}

\newcommand{}{\textit{Student answers orally; no code writing required.}}

\begin{document}

\begin{center}
    {\LARGE \textbf{DA3410 Viva Questions Set-1}}\\[4pt]
\end{center}

% Bank Q7
\begin{question}
What changes if \texttt{mean()} is replaced by \texttt{sum()} in this loss? 

\begin{lstlisting}
per_sample = (pred - target) ** 2
loss = per_sample.mean()
\end{lstlisting}
\end{question}

% Bank Q8
\begin{question}
Give the activation output and local derivative for the input below.
Use the convention that the ReLU derivative at zero is $0$.

\begin{center}
\begin{tabular}{lc}
\toprule
Activation and input & Required output \\
\midrule
ReLU, $x=3$ & $f(x)$ and $f'(x)$ \\
\bottomrule
\end{tabular}
\end{center}
\end{question}

% Bank Q10
\begin{question}
Consider the implementation below.

\begin{verbatim}
h = x @ W1 + b1
y = h @ W2 + b2
\end{verbatim}

Can the complete mapping from \texttt{x} to \texttt{y} always be replaced by
a single affine layer of the form

\begin{verbatim}
y = x @ W + b
\end{verbatim}

without changing the function? Explain briefly.
\end{question}

% Bank Q12
\begin{question}
The second line is replaced by \texttt{exp\_z = torch.exp(z)}. What changes? 

\begin{lstlisting}
z = torch.tensor([1000.0, 1001.0])
exp_z = torch.exp(z - z.max())
p = exp_z / exp_z.sum()
\end{lstlisting}
\end{question}

% Bank Q14
\begin{question}
Is the checkpointing code correct? If not, identify the problem and state the
fix conceptually.

\begin{verbatim}
for epoch in range(epochs):
    best_loss = float("inf")
    ...
    if val_loss < best_loss:
        best_loss = val_loss
        best_params = clone_params(params)
\end{verbatim}
\end{question}

% Bank Q19
\begin{question}
Consider the following manual autodiff forward pass:

\begin{verbatim}
a = multiply(x, w)
b = relu(a)
loss = square(b)
\end{verbatim}

The operations are recorded during the forward pass as

\begin{verbatim}
ops = [multiply_op, relu_op, square_op]
\end{verbatim}

Is the proposed backward implementation correct? If not, state the required
change.

\begin{verbatim}
for op in reversed(ops):
    op.backward()
\end{verbatim}
\end{question}

% Bank Q22
\begin{question}
What is the shape of \texttt{grad\_W}? 

\begin{lstlisting}
x = torch.randn(16, 10)
grad_z = torch.randn(16, 7)
grad_W = x.T @ grad_z
\end{lstlisting}
\end{question}

% Bank Q25
\begin{question}
Consider the following code for one training epoch:

\begin{verbatim}
updates = 0

for start in range(0, N, batch_size):
    xb = X[start:start + batch_size]
    update(params, xb)
    updates += 1

print(updates, len(xb))
\end{verbatim}

For the values below, what does the final \texttt{print} output?

\begin{center}
\begin{tabular}{cc}
\toprule
\texttt{N} & \texttt{batch\_size} \\
\midrule
300 & 300 \\
\bottomrule
\end{tabular}
\end{center}
\end{question}

% Bank Q29
\begin{question}
The validation loss improves, but the supposedly saved ``best'' parameters
keep changing during later training. Identify the bug. 

\begin{lstlisting}
if val_loss < best_loss:
    best_loss = val_loss
    best_params = params
\end{lstlisting}
\end{question}

% Bank Q32
\begin{question}
Suppose \texttt{group\_id[i]} identifies which related group sample
\texttt{i} belongs to. Samples with the same \texttt{group\_id} should not
appear in both training and validation.

Does the code guarantee this? If not, explain the problem.

\begin{verbatim}
groups = torch.unique(group_id)
groups = groups[torch.randperm(len(groups))]
cut = int(0.8 * len(groups))

train_groups = groups[:cut]
val_groups   = groups[cut:]
\end{verbatim}

The corresponding samples are then selected according to their group.
\end{question}


\end{document}
